% ---
% title: Low Yield Learning
% date: 2026-08
% category:
% nodes: reflections 
% links: 
% ---
\documentclass{article}
\usepackage[utf8]{inputenc}
\usepackage{amsmath, amssymb, amsthm}
\usepackage{geometry}
\usepackage{graphicx}
\usepackage{hyperref}

\geometry{a4paper, margin=1in}

\title{Low Yield Learning}
\date{August 2026}
\author{Daksh Mehta}

\begin{document}

IN PROGRESS

The vedas were recited word for word, passed through generations. Information is collated organised and available 24/7 at whatever level you need. Combined with an always expanding set of guides and tutorials. Knowledge workers.

Choosing what to learn 

Is there any point in learning things anymore?

Intention behind learning

Disproving the different learning theories and the book making it stick 

Will this come on the test? Is this high yield? 

Not caring about learning what is relevant or not. The New England Journal of Medicine regularly has sections about politics or health economics despite being a journal of medicine. 

Learning is a living, breathing form. Templates and guides can be helpful. But no amount of videos on best learning resources to learn x or y. No obsessing over templates for note taking. Pretending like taking notes in the Cornell format will give you 100 percent retention. Cornell format works through intention. 

Obsession about retention percentage. Trying to remember everything you've read. In my very anecdotal, non scientific wishy washy experience you probably retain somewhere between 10 to 70 percent of what you learn depending on the type of content (confidence interval large). Point being certain things can help boost retention but obsessing over those also delays the actual learning process. 

Adult learning - self motivated and not self motivated. How do you learn something which you don't want to? How do you learn things you do want to learn? Should you even view these things differently? 

Talking to people and teaching as a form of learning
	
\end{document}